\problemname{Book Reviews}
A book reviewer has read $N$ books that are to be reviewed. Each review should
end with the book being assigned a rating on a scale from $1$ to $M$. It can
be difficult to directly choose an absolute rating for each book, so
the reviewer thinks it is much easier to compare two books
at a time and describe which of them is better.

The reviewer has numbered the books with integers from $1$ to $N$ and now wants
to determine their ratings $a_1, a_2, \dots , a_N$. To do this, the
reviewer has made $R$ comparisons that describe the relation between $a_i$ and
$a_j$, for some books $i, j$.

The reviewer is satisfied with any rating assignment, as long as all requirements
from the comparisons are fulfilled. Help the reviewer find such a
rating assignment.

\section*{Input}

The first line consists of three integers, $N$ ($1 \leq N \leq 10^5$), 
$M$ ($1 \leq M \leq 10^5$),
$R$ ($1 \leq R \leq 5 \cdot 10^5$), the number of books, the highest possible
rating, and the number of comparisons.

Then follow $R$ lines with relations that must
be satisfied. Each such line has the format ``\texttt{<i> <relation> <j>}'',
which describes the relation between $a_i$ and $a_j$. $i$ and $j$ are integers between
$1$ and $N$, $i \neq j$. The relation is one of the strings `\texttt{<}', `\texttt{=}',
`\texttt{<=}', and this describes precisely that $a_i$ \texttt{<relation>} $a_j$ must hold.
No pair of books will be compared more than once.

\section*{Output}

Output a list of integers $a_1, a_2, \ldots , a_N$ such that all relations hold, and
all numbers are in the interval $[1, M]$. If there are multiple solutions, output any of them.
If it is impossible, output $-1$.

\section*{Scoring}
Your solution will be tested on a set of test groups.
To earn points for a group, you must pass all test cases in that group.

\noindent
\begin{tabular}{| l | l | l |}
  \hline
  \textbf{Group} & \textbf{Points} & \textbf{Constraints} \\ \hline
  $1$   & $30$       & Relations can only be \texttt{<} \\ \hline
  $2$   & $30$       & Relations can only be \texttt{<} or \texttt{=} \\ \hline
  $3$   & $40$       & No additional constraints. \\ \hline
\end{tabular}

\section*{Explanation of Sample 1}
In the first input example, \texttt{1 2 1 3 3} is a valid solution. This can be verified
by checking that all numbers are in the interval $[1, 3]$, and that the numbers satisfy the four relations $a_1 < a_2$,
$a_2 < a_4$, $a_3 < a_2$, and $a_2 < a_5$.
